\documentclass[12pt]{article}
\usepackage[utf8]{inputenc}
\usepackage[T1]{fontenc}
\usepackage{hyperref}
\title{Wulguru Language Resources for K-12 Curriculum Development}
\author{Compiled by Claude Sonnet 5.5 (claude-sonnet-5-5) for Andr\'as Kornai}
\date{}
\begin{document}
\maketitle

\section{Introduction}
Wulguru (ISO 639-3 qgu; Glottocode wulg1239) is an extinct Pama--Nyungan language of the Townsville area, Queensland. Documentation is fragmentary, drawing mainly on a late-19th-century journal by Charles Price, analysed in Mark Donohue's 2007 salvage study. The name is used as an umbrella for several related varieties. Austlang lists no known speakers.

\section{Primary Educational Resources}
\begin{itemize}
\item State Library of Queensland, \emph{Wulguru Everyday Words} (wordlist with suggested pronunciations): \url{https://content.slq.qld.gov.au/sites/default/files/Wulguru-Everyday-Words.pdf}
\end{itemize}

\section{Linguistic and Reference Materials}
\begin{itemize}
\item Wikipedia, Wulguru language: \url{https://en.wikipedia.org/wiki/Wulguru_language}
\item Glottolog record for Donohue 2007, \emph{Wulguru: A salvage study of the north-eastern Australian language from Townsville}: \url{https://glottolog.org/resource/reference/id/37337}
\item LINGUIST List issue (book announcement): \url{https://linguistlist.org/issues/17/3604}
\item AIATSIS Austlang entry Y135: \url{https://collection.aiatsis.gov.au/austlang/language/y135}
\end{itemize}

\section{Religious and Community Resources}
Revival work is reported to be supported by the North Queensland Regional Aboriginal Corporation Language Centre (website nqraclc.com.au was named in a result; its URL was not itself seen as a result and is not cited as a link).

\section{Digital and Audio Resources}
State Library of Queensland language-week posts:
\begin{itemize}
\item Wulgurukaba: \url{https://www.slq.qld.gov.au/blog/language-week-week-thirty-one-wulgurukaba}
\item Buluguyban: \url{https://www.slq.qld.gov.au/blog/language-week-week-nineteen-buluguyban}
\end{itemize}
Tsunoda's 1974 recordings of Buluguyban are reportedly held at AIATSIS; no direct link was found.

\section{Curriculum Development Context}
Material is limited to wordlists and a scholarly sketch. A realistic curriculum would be a revival-oriented vocabulary and cultural unit developed with the Wulgurukaba and Manbarra communities, not full instruction in the language.
\end{document}